\uuid{Ekyd}
\titre{Irrationnel}
\auteur{bodin ; hanani ; mzari}
\organisation{exo7}
\licence{CC BY-NC-SA 4.0}
\datecreate{2020-01-20}
\niveau{L1}
\module{Analyse}
\chapitre{Propriétés de R}
\sousChapitre{Les rationnels}
\difficulte{3}
\theme{Rationnels, irrationnel}
\exo7qcmid{lille-1.reels.006}

\contenu{
\begin{question}
\qclassification{120.01}
\qtags{motcle=irrationnel}
Je veux montrer que $\log 13$, est un nombre irrationnel. On rappelle que $\log 13$ est le réel tel que $10^{\log 13} = 13$. Quelle démarche puis-je adopter ?
\begin{answers}
    \bad{Par division je calcule l'écriture décimale de $\log 13$ et je montre qu'elle est périodique.}

    \bad{Je prouve par récurrence que $\log n$ est irrationnel pour $n\ge2$.}

    \good{Je suppose par l'absurde que $\log 13 = \frac pq$, avec $p,q \in \Nn^*$ et je cherche une contradiction après avoir écrit $13^q = 10^p$.}

    \bad{Il est faux que $\log 3$ soit un nombre irrationnel.}

\end{answers}
\begin{explanations}
On raisonne par l'absurde en écrivant $\log 13 = \frac pq$, où $p,q$ sont des entiers strictement positifs. On en déduit que $13^q = 10^p$. Comme $13$ et $10$ sont premiers entre eux, alors on obtient $p=q=0$ et donc une contradiction.
\end{explanations}
\end{question}
}
